%==============================================================================
%  MATH 231 -- Linear Algebra I (Queens College, Fall 2026)
%  MATRIX UNIT, PART 7 -- The Determinant in General
%
%  Section numbering runs CONTINUOUSLY across the parts of the matrix unit;
%  see Part 2's header for the full map.  This part is Sections 35-39, hence
%  \setcounter{section}{34}.
%
%  STATUS: complete.  Sections 35-39.  Written from the plan in
%  plan/MATH231_remaining_units_plan.tex (2026-10-25).
%
%  AUDIENCE: distributed to students.  No instructor-facing apparatus.
%
%  WHAT THE INSTRUCTOR SPECIFIED, AND WHERE IT IS.
%   - Generalizing the determinant: Laplace expansion, cofactor expansions,
%     minors (Sec. 35-36), principal minors (Sec. 37), general properties
%     (Sec. 38), determinants of orthogonal, diagonal and triangular matrices
%     with straightforward proofs (Sec. 39).
%
%  DECISIONS CARRIED FROM THE PLAN.
%
%   D6. WHAT IS PROVED.  The determinant is DEFINED by expansion along the
%       first row (Def. 36.1).  STATED without proof, because a proof needs
%       permutations: that expansion along any row or column gives the same
%       number (Thm 36.2); that a row swap flips the sign (D5); the product
%       rule for n > 2 (D9).  PROVED from those: det I = 1 (D1); det A' =
%       det A (D2, by induction using Thm 36.2); zero row (D3); scaling a row
%       and det(aA) = a^n det A (D4); two equal rows (D6, from D5); linearity
%       in a row (D7); row replacement (D8); dependent columns give 0 (half of
%       D10); triangular, diagonal, orthogonal (Sec. 39).  The other half of
%       D10 (rank n => det != 0) is completed in Sec. 42.6 by elimination.
%
%   The property list is labelled D1-D10 INSIDE THE DOCUMENT (determinant
%   laws), in the style of Part 1's M1-M6 and R1-R4.  These are unrelated to
%   the plan's decision numbers D1-D12 cited in these header comments.
%
%   PRINCIPAL MINORS (Sec. 37) close with Sylvester's criterion for a
%   symmetric 2x2, proved by completing the square, and a pointer to the
%   eigen unit and the Hessian.  The n x n criterion is stated only.
%
%  Compile with pdflatex or tectonic (standard TeX Live packages only).
%==============================================================================
\documentclass[11pt]{article}

\usepackage[margin=1in]{geometry}
\usepackage{amsmath,amssymb,amsthm}
\usepackage[dvipsnames]{xcolor}
\usepackage{enumitem}
\usepackage{booktabs}
\usepackage{array}
\usepackage[hidelinks]{hyperref}
\usepackage{titlesec}
\usepackage{needspace}

%------------------------------------------------------------------ style ----
\titleformat{\section}{\large\bfseries\sffamily}{\thesection.}{0.6em}{}
\titleformat{\subsection}{\normalsize\bfseries\sffamily}{\thesubsection.}{0.5em}{}
\titlespacing*{\section}{0pt}{14pt}{6pt}
\setlength{\emergencystretch}{3em}

\theoremstyle{plain}
\newtheorem{thm}{Theorem}[section]
\newtheorem{lem}[thm]{Lemma}
\newtheorem{cor}[thm]{Corollary}

\theoremstyle{definition}
\newtheorem{defn}{Definition}[section]
\newtheorem{ex}{Example}[section]
\theoremstyle{remark}
\newtheorem*{rmk}{Remark}
\newtheorem*{warn}{A common confusion}

\newcommand{\lookahead}{\par\smallskip\noindent
  \textbf{\sffamily\textcolor{RoyalBlue}{Looking ahead.}}\ \ignorespaces}
\newcommand{\lookback}[1][the vector unit]{\par\smallskip\noindent
  \textbf{\sffamily\textcolor{ForestGreen}{From #1.}}\ \ignorespaces}

%--------------------------------------------------------------- notation ----
\newcommand{\R}{\mathbb{R}}
\newcommand{\vc}[1]{\mathbf{#1}}
\newcommand{\mt}[1]{\mathbf{#1}}
\newcommand{\norm}[1]{\lVert #1\rVert}
\newcommand{\T}{^{\mathsf{T}}}
\newcommand{\nullsp}[1]{\operatorname{null}[#1]}
\DeclareMathOperator{\spanop}{span}
\DeclareMathOperator{\diagop}{diag}
\DeclareMathOperator{\rank}{rank}

\title{\vspace{-1.2cm}\textbf{\sffamily MATH 231 -- Linear Algebra I}\\[4pt]
       {\large\sffamily The Matrix Unit, Part 7 \textbullet\ The Determinant
        in General}\\[2pt]
       {\normalsize\sffamily Kiely Hall 326 \textbullet\
        Instructor: Kevin Bedoya}}
\author{}
\date{}

\begin{document}
\setcounter{section}{34}  % this part opens at Section 35; see the header note
\maketitle
\vspace{-1.6cm}

\noindent\textbf{About this document.} This is Part 7 of the matrix unit,
Sections 35--39. Section numbers run continuously through the unit (Part 2's
opening table lists them all).

\medskip
\noindent\textbf{Assumed.} The $2\times2$ determinant: its definition (\S14),
its properties P1--P11 (\S22), and its meaning as an area (\S23). Rank and
null space (\S12, \S31--\S32). Orthogonal matrices (\S20) and the rotations of
$\R^{3}$ (\S25).

\medskip
\noindent\textbf{What you should be able to do afterwards.} Compute minors and
cofactors; compute the determinant of a $3\times3$ or $4\times4$ matrix by
cofactor expansion, choosing a good row or column; find principal and leading
principal minors; use the properties of the determinant to compute it without
expanding; and prove that the determinant of a triangular matrix is the product
of its diagonal and that an orthogonal matrix has determinant $\pm1$.

\bigskip
\noindent For a $2\times2$ matrix the number $ad-bc$ did three jobs. It measured
how the matrix scales area, its sign recorded whether the plane was flipped, and
it was nonzero exactly when the matrix had full rank. We want a number with the
same three powers for an $n\times n$ matrix: a volume factor in $\R^{3}$, and the
same test of rank in every dimension. It is built recursively, out of
determinants one size smaller, with $ad-bc$ as the first step.

%==============================================================================
\section{Minors and cofactors}
%==============================================================================

%------------------------------------------------------------------------------
\subsection{Minors}
%------------------------------------------------------------------------------
\begin{defn}[Minor]
Let $\mt{A}\in\R^{n\times n}$ with $n\ge2$. Deleting row $i$ and column $j$ of
$\mt{A}$ leaves an $(n-1)\times(n-1)$ matrix, written $\mt{A}_{ij}$. Its
determinant
\[
  M_{ij}=\det(\mt{A}_{ij})
\]
is the \textbf{$(i,j)$ minor} of $\mt{A}$.
\end{defn}

\noindent The definition is recursive: a minor of an $n\times n$ matrix is the
determinant of a matrix one size smaller. For a $3\times3$ matrix the minors are
$2\times2$ determinants, which we know how to compute.

\begin{ex}
\label{ex:minors}
For
$\mt{A}=\begin{bmatrix}1&2&3\\4&5&6\\7&8&10\end{bmatrix}$, the three minors along
the first row are
\[
  M_{11}=\det\begin{bmatrix}5&6\\8&10\end{bmatrix}=2,\qquad
  M_{12}=\det\begin{bmatrix}4&6\\7&10\end{bmatrix}=-2,\qquad
  M_{13}=\det\begin{bmatrix}4&5\\7&8\end{bmatrix}=-3 .
\]
For $M_{12}$, row $1$ and column $2$ were crossed out, leaving the four entries
$4,6,7,10$.
\end{ex}

%------------------------------------------------------------------------------
\subsection{Cofactors}
%------------------------------------------------------------------------------
\begin{defn}[Cofactor]
The \textbf{$(i,j)$ cofactor} of $\mt{A}$ is the minor with a sign attached:
\[
  C_{ij}=(-1)^{i+j}M_{ij}.
\]
\end{defn}

\noindent The sign $(-1)^{i+j}$ is $+$ when $i+j$ is even and $-$ when it is
odd, which makes a checkerboard starting with $+$ in the top-left corner:
\[
  \begin{bmatrix}+&-&+&\cdots\\-&+&-&\cdots\\+&-&+&\cdots\\\vdots&\vdots&\vdots&\ddots\end{bmatrix}.
\]
In Example~\ref{ex:minors}, $C_{11}=+2$, $C_{12}=-(-2)=2$, $C_{13}=+(-3)=-3$.

%------------------------------------------------------------------------------
\subsection{The \texorpdfstring{$2\times2$}{2x2} case, seen again}
%------------------------------------------------------------------------------
Agree that the determinant of a $1\times1$ matrix $[\,a\,]$ is the number $a$
itself. Then for $\mt{A}=\begin{bmatrix}a&b\\c&d\end{bmatrix}$, deleting row 1 and
column 1 leaves $[\,d\,]$, and deleting row 1 and column 2 leaves $[\,c\,]$, so
$C_{11}=d$ and $C_{12}=-c$, and
\[
  aC_{11}+bC_{12}=ad-bc=\det\mt{A}.
\]
The $2\times2$ formula already has the shape ``each entry of the first row,
times its cofactor, added up''. That shape is the general definition.

%==============================================================================
\section{The Laplace expansion}
%==============================================================================

%------------------------------------------------------------------------------
\subsection{The definition}
%------------------------------------------------------------------------------
\begin{defn}[Determinant]
For $\mt{A}\in\R^{n\times n}$: if $n=1$, $\det[\,a\,]=a$. If $n\ge2$,
\[
  \boxed{\;\det\mt{A}=\sum_{j=1}^{n}a_{1j}C_{1j}
  =a_{11}C_{11}+a_{12}C_{12}+\cdots+a_{1n}C_{1n}.\;}
\]
This is the \textbf{Laplace expansion}, or \textbf{cofactor expansion}, along the
first row.
\end{defn}

\noindent Each cofactor is a determinant of size $n-1$, computed by the same rule,
until the sizes reach $2\times2$ or $1\times1$.

\begin{ex}
For the matrix of Example~\ref{ex:minors},
\[
  \det\mt{A}=1\cdot C_{11}+2\cdot C_{12}+3\cdot C_{13}=1(2)+2(2)+3(-3)=-3 .
\]
\end{ex}

%------------------------------------------------------------------------------
\subsection{The \texorpdfstring{$3\times3$}{3x3} formula}
%------------------------------------------------------------------------------
Written out for a general $3\times3$ matrix, the definition gives
\[
  \det\begin{bmatrix}a_{11}&a_{12}&a_{13}\\a_{21}&a_{22}&a_{23}\\a_{31}&a_{32}&a_{33}\end{bmatrix}
  =a_{11}(a_{22}a_{33}-a_{23}a_{32})
   -a_{12}(a_{21}a_{33}-a_{23}a_{31})
   +a_{13}(a_{21}a_{32}-a_{22}a_{31}),
\]
a sum of six products. Each product takes exactly one entry from each row and
one from each column. That is true in every size, and it is the starting point of
the closed formula (with permutations) that more advanced courses use.

%------------------------------------------------------------------------------
\subsection{Any row, any column}
%------------------------------------------------------------------------------
The first row has no special status.

\begin{thm}[Expansion along any row or column]
\label{thm:anyrow}
For $\mt{A}\in\R^{n\times n}$ and any fixed $i$ and $j$,
\[
  \det\mt{A}=\sum_{k=1}^{n}a_{ik}C_{ik}
  \quad(\text{along row } i),
  \qquad
  \det\mt{A}=\sum_{k=1}^{n}a_{kj}C_{kj}
  \quad(\text{down column } j).
\]
\end{thm}

\noindent All $2n$ expansions give the same number. The proof needs the closed
formula just mentioned and is beyond this course; we take the theorem as given.
It can at least be checked.

\begin{ex}
Expand the same $\mt{A}$ down its first column. The cofactors needed are
$C_{11}=2$,
$C_{21}=-\det\begin{bmatrix}2&3\\8&10\end{bmatrix}=-(20-24)=4$ and
$C_{31}=+\det\begin{bmatrix}2&3\\5&6\end{bmatrix}=12-15=-3$, so
\[
  \det\mt{A}=1(2)+4(4)+7(-3)=2+16-21=-3 . \quad\checkmark
\]
\end{ex}

%------------------------------------------------------------------------------
\subsection{Choosing wisely}
%------------------------------------------------------------------------------
Every zero entry in the chosen row or column kills a whole term, cofactor and
all. So expand along the row or column with the most zeros.

\begin{ex}
$\mt{B}=\begin{bmatrix}2&0&1\\3&0&4\\1&5&2\end{bmatrix}$. Column 2 has two zeros,
so only $b_{32}=5$ contributes:
\[
  \det\mt{B}=5\,C_{32}=5\cdot(-1)^{3+2}\det\begin{bmatrix}2&1\\3&4\end{bmatrix}
  =-5(8-3)=-25 .
\]
Along the first row instead:
$2\det\begin{bmatrix}0&4\\5&2\end{bmatrix}-0+1\det\begin{bmatrix}3&0\\1&5\end{bmatrix}
=2(-20)+15=-25$. Same answer, more work.
\end{ex}

\begin{ex}
A $4\times4$ example, where the choice matters more:
\[
  \mt{C}=\begin{bmatrix}1&0&0&2\\0&3&0&0\\0&0&2&0\\1&0&0&1\end{bmatrix}.
\]
Row 2 has a single nonzero entry, $c_{22}=3$, with sign $(-1)^{2+2}=+$:
\[
  \det\mt{C}=3\det\begin{bmatrix}1&0&2\\0&2&0\\1&0&1\end{bmatrix}.
\]
The $3\times3$ determinant has a single nonzero entry in its middle row, $2$,
again with sign $+$: it equals $2\det\begin{bmatrix}1&2\\1&1\end{bmatrix}=2(1-2)=-2$.
So $\det\mt{C}=3(-2)=-6$.
\end{ex}

%------------------------------------------------------------------------------
\subsection{The cost}
%------------------------------------------------------------------------------
Expanding an $n\times n$ determinant creates $n$ determinants of size $n-1$,
each of which creates $n-1$ of size $n-2$, and so on: about
$n\cdot(n-1)\cdots2=n!$ products in all. For $n=10$ that is about $3.6$ million;
for $n=20$, about $2.4\times10^{18}$, more than a fast computer can do in a
year. Cofactor expansion is a definition and a tool for small or sparse
matrices, not an algorithm for large ones. Part 8 computes determinants by
elimination, with about $\tfrac23n^{3}$ operations: $5{,}000$-odd for $n=20$.

%------------------------------------------------------------------------------
\subsection{Volume}
%------------------------------------------------------------------------------
For a $3\times3$ matrix, $|\det\mt{A}|$ is the volume of the parallelepiped (the
slanted box) with edges the three columns, and every volume is multiplied by
$|\det\mt{A}|$ when $\mt{A}$ is applied. This is \S23 one dimension up; we state
it without proof. The sign records orientation: positive if the columns, in
order, are arranged like $\vc{e}_1,\vc{e}_2,\vc{e}_3$ (a right-handed
arrangement), negative if like their mirror image. For a box with edges along
the axes,
$\det\diagop(a,b,c)=abc$, the familiar length times width times height
(\S39.2).

%==============================================================================
\section{Principal minors}
%==============================================================================

%------------------------------------------------------------------------------
\subsection{The definition}
%------------------------------------------------------------------------------
A minor deletes one row and one column. A \emph{principal} minor keeps the rows
and columns in matching sets.

\begin{defn}[Principal and leading principal minors]
Let $\mt{A}\in\R^{n\times n}$ and choose a set $S$ of indices from
$\{1,\dots,n\}$. Keeping the rows \emph{and} the columns numbered by $S$ gives a
\textbf{principal submatrix}; its determinant is a \textbf{principal minor}.
When $S=\{1,2,\dots,k\}$, the submatrix is the top-left $k\times k$ corner, and
its determinant $\Delta_k$ is the \textbf{$k$-th leading principal minor}.
\end{defn}

\noindent So $\Delta_1=a_{11}$, $\Delta_2=\det\begin{bmatrix}a_{11}&a_{12}\\a_{21}&a_{22}\end{bmatrix}$,
and so on up to $\Delta_n=\det\mt{A}$. A principal submatrix always has its own
diagonal on the diagonal of $\mt{A}$, which is the point of keeping the same
indices.

%------------------------------------------------------------------------------
\subsection{An example}
%------------------------------------------------------------------------------
\begin{ex}
\label{ex:principal}
$\mt{A}=\begin{bmatrix}2&1&0\\1&3&1\\0&1&4\end{bmatrix}$.
\begin{itemize}[leftmargin=1.4em,itemsep=3pt]
  \item Size $1$ ($S=\{1\},\{2\},\{3\}$): the diagonal entries $2,3,4$.
  \item Size $2$: $S=\{1,2\}$ gives $\det\begin{bmatrix}2&1\\1&3\end{bmatrix}=5$;
        $S=\{1,3\}$ gives $\det\begin{bmatrix}2&0\\0&4\end{bmatrix}=8$;
        $S=\{2,3\}$ gives $\det\begin{bmatrix}3&1\\1&4\end{bmatrix}=11$.
  \item Size $3$: $\det\mt{A}=2(12-1)-1(4-0)+0=18$.
\end{itemize}
The leading principal minors are $\Delta_1=2$, $\Delta_2=5$, $\Delta_3=18$.
\end{ex}

\noindent Note that the size-$1$ principal minors add up to the trace (\S9). The
size-$2$ ones and the size-$n$ one (the determinant) also have a meaning, which
the next unit uncovers.

%------------------------------------------------------------------------------
\subsection{What they are for}
%------------------------------------------------------------------------------
A symmetric matrix $\mt{A}$ is called \textbf{positive definite} if
$\vc{x}\T\mt{A}\vc{x}>0$ for every $\vc{x}\neq\vc{0}$. (The expression
$\vc{x}\T\mt{A}\vc{x}$ appeared in \S20.) Such matrices are the matrix version of
a positive number. They are the Hessians of functions shaped like a bowl, and
they appear in least squares in Part 9. The leading principal minors detect them.

\begin{thm}[Sylvester's criterion, $2\times2$]
A symmetric $\mt{A}=\begin{bmatrix}a&b\\b&c\end{bmatrix}$ is positive definite if
and only if $\Delta_1=a>0$ and $\Delta_2=ac-b^{2}>0$.
\end{thm}

\begin{proof}
Here $\vc{x}\T\mt{A}\vc{x}=ax_1^{2}+2bx_1x_2+cx_2^{2}$. If $a>0$, complete the
square:
\[
  ax_1^{2}+2bx_1x_2+cx_2^{2}
  =a\Bigl(x_1+\frac{b}{a}x_2\Bigr)^{2}+\frac{ac-b^{2}}{a}\,x_2^{2}.
\]
If also $ac-b^{2}>0$, both terms are $\ge0$, and both are $0$ only if $x_2=0$ and
then $x_1=0$. Conversely, if $\mt{A}$ is positive definite then $\vc{x}=\vc{e}_1$
gives $a>0$, and $\vc{x}=[\,-b,a\,]$ makes the first term vanish and gives
$a(ac-b^{2})>0$, so $ac-b^{2}>0$.
\end{proof}

\noindent For $n\times n$ the same holds: a symmetric matrix is positive
definite exactly when all $n$ leading principal minors are positive. We state this
without proof. Example~\ref{ex:principal}'s matrix ($2,5,18$) is positive
definite. Applied to a Hessian (the vector calculus unit, \S15), this is the
multivariable second-derivative test for a minimum.

%==============================================================================
\section{Properties of the determinant}
%==============================================================================
These are the laws that make the determinant computable without expanding, and
they generalize \S22 property by property. They are labelled D1--D10, the way
Part 1 labelled the laws of multiplication M1--M6. Throughout,
$\mt{A},\mt{B}\in\R^{n\times n}$.

%------------------------------------------------------------------------------
\subsection{The laws}
%------------------------------------------------------------------------------
\begin{center}
\renewcommand{\arraystretch}{1.4}
\begin{tabular}{@{}l p{8.4cm} l@{}}
\toprule
 & \textbf{Law} & \textbf{Status} \\
\midrule
\textbf{D1} & $\det\mt{I}_n=1$ & proved \\
\textbf{D2} & $\det\mt{A}\T=\det\mt{A}$ & proved, from Theorem~\ref{thm:anyrow} \\
\textbf{D3} & a zero row gives $\det\mt{A}=0$ & proved \\
\textbf{D4} & scaling one row by $\alpha$ scales $\det$ by $\alpha$; so $\det(\alpha\mt{A})=\alpha^{n}\det\mt{A}$ & proved \\
\textbf{D5} & swapping two rows changes the sign & stated \\
\textbf{D6} & two equal rows give $\det\mt{A}=0$ & proved, from D5 \\
\textbf{D7} & $\det$ is linear in each row separately & proved \\
\textbf{D8} & adding a multiple of one row to another leaves $\det$ unchanged & proved, from D6--D7 \\
\textbf{D9} & $\det(\mt{A}\mt{B})=\det\mt{A}\,\det\mt{B}$ & stated ($n=2$ in \S22) \\
\textbf{D10} & $\det\mt{A}\neq0\iff\rank\mt{A}=n\iff\nullsp{\mt{A}}=\{\vc{0}\}$ & half here, half in \S42.6 \\
\bottomrule
\end{tabular}
\end{center}

\noindent By D2, every law about rows is also a law about columns.

%------------------------------------------------------------------------------
\subsection{The proofs}
%------------------------------------------------------------------------------
\textbf{D1.} Expand $\mt{I}_n$ along its first row: the only nonzero term is
$1\cdot\det\mt{I}_{n-1}$. Repeat down to $\det\mt{I}_1=1$.

\textbf{D2.} By induction on $n$; for $n=1$ there is nothing to prove. Expand
$\mt{A}\T$ along its first row. Its entries are $a_{11},a_{21},\dots,a_{n1}$,
the first column of $\mt{A}$. Deleting row $1$ and column $j$ of $\mt{A}\T$
leaves the transpose of $\mt{A}_{j1}$, whose determinant equals
$\det\mt{A}_{j1}$ by the induction hypothesis. So the expansion of $\mt{A}\T$
along row $1$ is term by term the expansion of $\mt{A}$ down column $1$, which is
$\det\mt{A}$ by Theorem~\ref{thm:anyrow}.

\textbf{D3.} Expand along the zero row: every term has a factor $0$.

\textbf{D4.} Expand along the scaled row. Every term $a_{ik}C_{ik}$ becomes
$(\alpha a_{ik})C_{ik}$, while the cofactors $C_{ik}$ are unchanged, because
they never use row $i$. Scaling all $n$ rows multiplies by $\alpha$ $n$ times.
So $\det(2\mt{I}_3)=8$: doubling every length in space multiplies every volume
by $8$.

\textbf{D5} is stated. On Example~\ref{ex:minors}'s matrix, swapping rows $1$
and $2$ gives
$\det\begin{bmatrix}4&5&6\\1&2&3\\7&8&10\end{bmatrix}=4(-4)-5(-11)+6(-6)=3=-(-3)$.

\textbf{D6.} Swapping the two equal rows leaves $\mt{A}$ unchanged but, by D5,
changes the sign of the determinant. So $\det\mt{A}=-\det\mt{A}$, and
$\det\mt{A}=0$.

\textbf{D7.} If row $i$ is $\alpha\vc{u}+\beta\vc{v}$, expand along row $i$: each
term is $(\alpha u_k+\beta v_k)C_{ik}$, and the cofactors do not involve row $i$.
Splitting the sum gives $\alpha$ times the determinant with row $i$ equal to
$\vc{u}$ plus $\beta$ times the one with row $i$ equal to $\vc{v}$.

\textbf{D8.} Replace row $i$ by $\vc{r}_i+k\vc{r}_j$ with $j\neq i$. By D7 the new
determinant is $\det\mt{A}$ plus $k$ times the determinant of the matrix whose
rows $i$ and $j$ are both $\vc{r}_j$, which is $0$ by D6.
\[
  \det\begin{bmatrix}1&2&3\\0&-3&-6\\7&8&10\end{bmatrix}
  =1(-30+48)-2(0+42)+3(0+21)=-3,
\]
the determinant of Example~\ref{ex:minors}'s matrix after subtracting $4$ times
row $1$ from row $2$. Unchanged.

\textbf{D10, first half.} \emph{If the columns are dependent, $\det\mt{A}=0$.}
Suppose column $k$ is a combination of the others,
$\vc{a}_k=\sum_{j\neq k}c_j\vc{a}_j$. By D7 for columns, $\det\mt{A}$ is
$\sum_{j\neq k}c_j$ times the determinant of the matrix with $\vc{a}_j$ written in
column $k$. Each of those has two equal columns, so is $0$ by D6 for columns.
Hence $\rank\mt{A}<n$ forces $\det\mt{A}=0$. The converse, that $\rank\mt{A}=n$
forces $\det\mt{A}\neq0$, comes out of elimination in \S42.6.

%------------------------------------------------------------------------------
\subsection{What the row operations do}
%------------------------------------------------------------------------------
Three of the laws describe what happens to the determinant under the three
simplest ways of changing rows:
\begin{center}
\renewcommand{\arraystretch}{1.35}
\begin{tabular}{@{}l l@{}}
\toprule
\textbf{Operation on the rows} & \textbf{Effect on $\det$} \\
\midrule
swap two rows & multiply by $-1$ (D5) \\
multiply one row by $\alpha\neq0$ & multiply by $\alpha$ (D4) \\
add a multiple of one row to another & no change (D8) \\
\bottomrule
\end{tabular}
\end{center}
These are exactly the three operations of Gaussian elimination in Part 8. That
is why elimination computes determinants cheaply.

\begin{warn}[No law for sums]
There is a law for products, D9, and none for sums: in general
$\det(\mt{A}+\mt{B})\neq\det\mt{A}+\det\mt{B}$. Already
$\det(\mt{I}_3+\mt{I}_3)=8$ while $\det\mt{I}_3+\det\mt{I}_3=2$.
\end{warn}

%==============================================================================
\section{Determinants of special matrices}
%==============================================================================

%------------------------------------------------------------------------------
\subsection{Triangular matrices}
%------------------------------------------------------------------------------
\begin{thm}
\label{thm:triangular}
If $\mt{A}\in\R^{n\times n}$ is triangular (upper or lower), its determinant is
the product of its diagonal entries:
\[
  \det\mt{A}=a_{11}a_{22}\cdots a_{nn}.
\]
\end{thm}

\begin{proof}
Lower triangular first, by induction on $n$; the case $n=1$ is the definition.
Row $1$ of a lower triangular matrix is $[\,a_{11},0,\dots,0\,]$, so expanding
along it leaves one term:
\[
  \det\mt{A}=a_{11}C_{11}=a_{11}\det\mt{A}_{11}.
\]
Deleting row $1$ and column $1$ of a lower triangular matrix leaves a lower
triangular matrix of size $n-1$ with diagonal $a_{22},\dots,a_{nn}$, whose
determinant is $a_{22}\cdots a_{nn}$ by the induction hypothesis. For an upper
triangular $\mt{A}$, the transpose is lower triangular with the same diagonal, and
D2 finishes the proof.
\end{proof}

\begin{ex}
$\det\begin{bmatrix}2&5&-1\\0&3&7\\0&0&4\end{bmatrix}=2\cdot3\cdot4=24$. The entries
above the diagonal play no role at all.
\end{ex}

\noindent In particular, a triangular matrix has determinant $0$ exactly when
some diagonal entry is $0$. In Part 8 that becomes the condition under which a
triangular system can be solved by back substitution.

%------------------------------------------------------------------------------
\subsection{Diagonal matrices}
%------------------------------------------------------------------------------
A diagonal matrix is triangular, so $\det\diagop(d_1,\dots,d_n)=d_1d_2\cdots d_n$.
Geometrically, $\diagop(d_1,\dots,d_n)$ stretches the $k$-th axis by $d_k$, and a
unit cube becomes a box with edges $|d_1|,\dots,|d_n|$: the volume factor is the
product. With every $d_k=1$ this is D1 again.

%------------------------------------------------------------------------------
\subsection{Orthogonal matrices}
%------------------------------------------------------------------------------
\begin{thm}
If $\mt{Q}\in\R^{n\times n}$ is orthogonal, then $\det\mt{Q}=\pm1$.
\end{thm}

\begin{proof}
From $\mt{Q}\T\mt{Q}=\mt{I}$, by D9, D2 and D1,
$(\det\mt{Q})^{2}=\det\mt{Q}\T\det\mt{Q}=\det(\mt{Q}\T\mt{Q})=\det\mt{I}=1$.
\end{proof}

\noindent An orthogonal matrix preserves lengths and angles (\S20), so it
preserves volume, and the theorem confirms the factor is $1$; the sign records
whether space is turned over. The rotations of \S25 have determinant $+1$:
\[
  \det\mt{R}_z(\theta)=\det\begin{bmatrix}\cos\theta&-\sin\theta&0\\\sin\theta&\cos\theta&0\\0&0&1\end{bmatrix}
  =1\cdot\det\mt{R}_\theta=1,
\]
expanding along the third row. The same works for $\mt{R}_x$ (along the first
row) and $\mt{R}_y$ (along the second). The swap of two coordinates,
$\begin{bmatrix}0&1&0\\1&0&0\\0&0&1\end{bmatrix}$, is $\mt{I}_3$ with two rows
exchanged, so by D5 and D1 its determinant is $-1$: a reflection across the plane
$x=y$.

\begin{warn}[The converse is false]
$\det\mt{A}=\pm1$ does not make $\mt{A}$ orthogonal. The shear
$\begin{bmatrix}1&1\\0&1\end{bmatrix}$ has determinant $1$ and changes lengths
(\S21.4).
\end{warn}

%------------------------------------------------------------------------------
\subsection{Projections}
%------------------------------------------------------------------------------
\begin{thm}
If $\mt{P}^{2}=\mt{P}$ and $\mt{P}\neq\mt{I}$, then $\det\mt{P}=0$.
\end{thm}

\begin{proof}
Since $\mt{P}\neq\mt{I}$, there is a vector $\vc{x}$ with
$\vc{z}=(\mt{I}-\mt{P})\vc{x}\neq\vc{0}$. Then
$\mt{P}\vc{z}=(\mt{P}-\mt{P}^{2})\vc{x}=\vc{0}$, so $\nullsp{\mt{P}}$ contains
$\vc{z}\neq\vc{0}$. By Theorem~31.2 of Part 6 the columns of $\mt{P}$ are
dependent, and by the first half of D10, $\det\mt{P}=0$.
\end{proof}

\noindent So a projection that is not the identity always loses volume
entirely: it flattens space onto a lower-dimensional subspace.

%------------------------------------------------------------------------------
\subsection{Practice}
%------------------------------------------------------------------------------
\begin{enumerate}[label=\textbf{P\arabic*.},leftmargin=2.6em,itemsep=5pt]
  \item Compute $\det\begin{bmatrix}2&-1&0\\1&3&2\\0&1&4\end{bmatrix}$ along the
        first row, and again down the first column.
  \item Compute the determinant of the upper triangular matrix with diagonal
        $2,-1,3,\frac12$ (and anything above the diagonal).
  \item $\mt{A}$ is $3\times3$ with $\det\mt{A}=5$. Find $\det(2\mt{A})$,
        $\det(-\mt{A})$, $\det(\mt{A}\T\mt{A})$, and the determinant of the matrix
        obtained by swapping two rows of $\mt{A}$.
  \item Find the leading principal minors of
        $\begin{bmatrix}4&2&0\\2&3&1\\0&1&2\end{bmatrix}$. Is it positive definite?
  \item Explain without computing why
        $\det\begin{bmatrix}1&2&3\\2&4&6\\1&0&1\end{bmatrix}=0$.
  \item For which $k$ is
        $\det\begin{bmatrix}1&1&1\\1&2&k\\1&4&k^{2}\end{bmatrix}=0$? Explain the
        answer by looking at the columns.
\end{enumerate}

\subsubsection*{Answers}

\noindent\textbf{P1.} Along row 1: $2(12-2)-(-1)(4-0)+0=20+4=24$. Down column 1:
$2(12-2)-1\cdot\det\begin{bmatrix}-1&0\\1&4\end{bmatrix}+0=20-(-4)=24$.
\checkmark

\smallskip
\noindent\textbf{P2.} $2\cdot(-1)\cdot3\cdot\frac12=-3$.

\smallskip
\noindent\textbf{P3.} $2^{3}\cdot5=40$; $(-1)^{3}\cdot5=-5$; $5\cdot5=25$;
$-5$.

\smallskip
\noindent\textbf{P4.} $\Delta_1=4$, $\Delta_2=12-4=8$,
$\Delta_3=4(6-1)-2(4-0)+0=12$. All positive, so it is positive definite.

\smallskip
\noindent\textbf{P5.} Row 2 is twice row 1. By D8, subtracting twice row 1 from
row 2 does not change the determinant and produces a zero row, so by D3 the
determinant is $0$.

\smallskip
\noindent\textbf{P6.} Along the first row:
$(2k^{2}-4k)-(k^{2}-k)+(4-2)=k^{2}-3k+2=(k-1)(k-2)$, so $k=1$ or $k=2$. The
columns are $[\,1,x,x^{2}\,]$ for $x=1,2,k$; when $k$ equals $1$ or $2$, two
columns coincide.

\lookahead Part 8 finally answers the question Part 1 left open: how to solve
$\mt{A}\vc{x}=\vc{b}$. The method uses the three row operations of \S38.3 to
make the matrix triangular, and as a by-product it computes the determinant
with a few thousand operations instead of $n!$.

\end{document}
